%%%%%%%%%%%%%%%%%%%%%%%%%%%%%%%%%%%%%%%%%%%%%%%%%%%%%%%%%%%%%%%%%%%%%%%%%%%%%%%%
%THESES ORLEANAISES : STYLES POUR LES PAGES DE COUV et DE DOS obligatoires
%helene.jouguet@univ-orleans.fr / SCD

%Classe du document : book :
\documentclass[a4paper,12pt, twoside]{book}
%%%%%%%%%%%%%%%%%%%%%%%%%%%%%%%%%%%%%%%%%%%%%%%%%%%%%%%%%%%%%%%%%%%%%%%%%%%%%%%%
% Packages ---------------------------------------------------------------------
\usepackage[T1]{fontenc}			%codage de caractères T1
\usepackage[utf8]{inputenc}			%caractères accentués
\usepackage[francais]{babel}		%adaptation de LaTeX au français
\usepackage{xspace}					%Gestion des espaces
\usepackage{graphicx}				%inclure des graphiques
%\usepackage{vmargin}				%Definition des marges
%\setmarginsrb{2cm}{1cm}{2cm}{1cm}{0cm}{1cm}{1cm}{1cm}
%\usepackage[scaled]{uarial}   %arial pour les pages de couv et 4ème

%\usepackage{siunitx} % Pour les unités



%\usepackage[longnamesfirst]{natbib}

\usepackage[francais]{layout} %Incompatible avec vmargin
%\usepackage{fullpage}

%\usepackage[left=1.2in, right=1in, top=1in, bottom=1in]{geometry}
%\usepackage[left=3cm, right=2.5cm, top=2.5cm, bottom=2.5cm, bindingoffset=0.5cm]{geometry}
\usepackage[includeheadfoot, margin=2.6cm, bindingoffset=0.5cm]{geometry}
%\geometry{showframe=true}

\setlength{\marginparwidth}{60pt}
\setlength{\headheight}{28pt} % apparemment bien pour fancyhdr
\usepackage[textwidth=60pt,textsize=footnotesize]{todonotes}


\usepackage{fancyhdr}

% Autres packages non utilisés pour les couv.
\usepackage{lineno}					% Numéro de lignes
\usepackage{textcomp} 				% améliore l'affichage des symbol (degré °) TEST
\usepackage{amsmath}				% complement mathematique
\usepackage{amssymb}				% ajoute des symboles
%\usepackage{array}					% Outils supplémentaires pour les tableaux
%\usepackage{multicol}				% Plusieurs colonnes possibles
%\usepackage{indentfirst} 			% indenter le premier paragraphe après une nouvelle section
%\usepackage{latexsym}				% ajoute des symboles
%\usepackage{setspace}				% Pour les Interlignes
%\usepackage{algorithm}				% Pour les algorithmes
%\usepackage{algorithmic}			% Pour les algorithmes
\usepackage{makeidx}				% Pour l'index
%\usepackage{xspace}					% Permet les espace après macro newcommand
\usepackage{apalike}				% Formattage de la biblio
\usepackage{setspace}				% Interligne double

\usepackage[export]{adjustbox}		%allow to align images by the top
\usepackage{framed}					%allow to do the box arround the title
%\usepackage{pdfpages}
%\usepackage{quotchap}
\usepackage{booktabs}
\usepackage{titlesec} %Formattage des sction
\usepackage[colorlinks = true, allcolors=blue]{hyperref}


% Formattage headers footers ---------------------------------------------------
\fancyhf{}

% Définition des différents style de header
\fancypagestyle{plain}{ % Redéfinition des entête et pied pour page chapitre
\fancyhf{}
\fancyfoot[C]{\thepage}
\renewcommand{\headrulewidth}{0pt}
\renewcommand{\footrulewidth}{0pt}
}

\fancypagestyle{frontmatter}{% Définition des entête et pied de page pour frontmatter
  \renewcommand{\headrulewidth}{0pt}% No header rule
  \renewcommand{\footrulewidth}{0pt}% No footer rule
  \fancyhf{}% Clear header/footer
  \fancyfoot[C]{\thepage}%
}

\fancypagestyle{mainmatter}{% Définition des entête et pied de page pour frontmatter
  \fancyhf{}
  \fancyhead[LE]{\normalfont\nouppercase{\leftmark}}
  \fancyfoot[LE,RO]{\thepage} % partie gauche du pied de page
  \fancyhead[RO]{\normalfont\nouppercase{\rightmark}} % partie droite de l'en-tête
  \fancyfoot[RO]{\thepage}  % partie droite du pied de page
  \renewcommand{\headrulewidth}{0.4pt}
}

\fancypagestyle{backmatter}{% Définition des entête et pied de page pour frontmatter
  \renewcommand{\headrulewidth}{0pt}% No header rule
  \renewcommand{\footrulewidth}{0pt}% No footer rule
  \fancyhf{}% Clear header/footer
  \fancyfoot[C]{\thepage}%
}

\fancypagestyle{emptyPerso}{% Définition des entête et pied de page pour frontmatter
  \renewcommand{\headrulewidth}{0pt}% No header rule
  \renewcommand{\footrulewidth}{0pt}% No footer rule
  \fancyhf{}% Clear header/footer
}
% Redéfinition de chpatermark et sectionmark
\renewcommand{\chaptermark}[1]{\markboth{\chaptername\ \thechapter.\ #1}{}}
\renewcommand{\sectionmark}[1]{\markright{\thesection.\ #1}}

% Formattage des sections ------------------------------------------------------
%\titleformat{\chapter}{\Huge}{\thechapter}{1em}{}
\titleformat{\section}{\huge}{\thesection}{1em}{} %Les Sections

% Autres -----------------------------------------------------------------------
\makeindex								% "lancer" le log de l'index

%%%%%%%%%%%%%%%%%%%%%%%%%%%%%%%%%%%%%%%%%%%%%%%%%%%%%%%%%%%%%%%%%%%%%%%%%%%%%%%%
%%%%%%%%%%%%%%%%%%%%%%%%%%%%%%%%%%%%%%%%%%%%%%%%%%%%%%%%%%%%%%%%%%%%%%%%%%%%%%%%
%%%%%%%%%%%%%%%%%%%%%%%%%%%%%%%%%%%%%%%%%%%%%%%%%%%%%%%%%%%%%%%%%%%%%%%%%%%%%%%%
%%%%%%%%%%%%%%%%%%%%%%%%%%%%%%%%%%%%%%%%%%%%%%%%%%%%%%%%%%%%%%%%%%%%%%%%%%%%%%%%
\begin{document}

% Franciser certaines commandes
\renewcommand{\tablename}{Tableau}
%\renewcommand{\figurename}{Nom de figure}
\renewcommand{\listfigurename}{Liste des figures}%
%\renewcommand{\listtablename}{Nouveau nom}%

% Ajout de commandes utiles
\newcommand\plop{\textbf{(Réf needed)\xspace}}
\newcommand\COO{CO$_{2}$\xspace}
\newcommand\CHH{CH$_{4}$\xspace}

%%%%%%%%%%%%%%%%%%%%%%%%%%%%%%%%%%%%%%%%%%%%%%%%%%%%%%%%%%%%%%%%%%%%%%%%%%%%%%%%
% Frontmatter
%%%%%%%%%%%%%%%%%%%%%%%%%%%%%%%%%%%%%%%%%%%%%%%%%%%%%%%%%%%%%%%%%%%%%%%%%%%%%%%%
\frontmatter%
\pagestyle{emptyPerso}
%%%%%%%%%%%%%%%%%%%%%%%%%%%%%%%%%%%%%%%%%%%%%%%%%%%%%%%%%%%%%%%%%%%%%%%%%%%%%%%%
% 1e de couverture
%%%%%%%%%%%%%%%%%%%%%%%%%%%%%%%%%%%%%%%%%%%%%%%%%%%%%%%%%%%%%%%%%%%%%%%%%%%%%%%%
% set-up -----------------------------------------------------------------------
\pagenumbering{gobble} % stop page numbering
\newgeometry{left=2cm,top=1.5cm,right=1.5cm,bottom=1.5cm} % set up margin
\usefont{T1}{phv}{m}{n} % Set up font style
{\parindent0pt % disables indentation for all the text between { and }

% logos and university name ----------------------------------------------------
\includegraphics[width=0.2\textwidth, valign=c]{./images/logos/pucvl}
\hfill
{\LARGE\textbf{UNIVERSITÉ D'ORLÉANS}}
\hfill
\includegraphics[width=0.2\textwidth, valign=c]{./images/logos/univ}

% PhD Informations -------------------------------------------------------------
\begin{center}
	{\large\textit{\textbf{ÉCOLE DOCTORALE ÉNERGIE, MATÉRIAUX, SCIENCES DE LA TERRE ET DE L'UNIVERS}}}\\

	\vspace{1.1cm}

	{\large Laboratoire de Physique et Chimie de l'Environnement et de l'Espace} \\

	\vspace{1.1cm}

	{\Large\textbf{THÈSE}} présentée par :\\
	{\Large\textbf{Benoît D'ANGELO}}\\

	\vspace{1.1cm}

	\begin{tabular}{r l}
		soutenue le : & \textbf{[XX décembre 2015]} \\ [.8cm]
		pour obtenir le grade de : & \textbf{Docteur de l'universit\'e d'Orl\'eans}\\[.2cm]
		Discipline : &  \textbf{Sciences de la Terre et de l'Univers}\\
	\end{tabular}
\end{center}

%\vspace{0.2cm}
\vfill

% Title box --------------------------------------------------------------------
\begin{framed}
	\begin{minipage}{\dimexpr\textwidth-2\fboxrule-2\fboxsep}
	\centering
		\vspace{0.2cm}
			{\Large\textbf{Variabilité spatio-temporelle des émissions de GES dans une tourbière à Sphaignes : \\effets sur le bilan de carbone}\\ \vspace{0.3cm}
			}
		\vspace{0.2cm}
	\end{minipage}
\end{framed}

} % end of the "no indentation" parametrization

\vfill

% Encadrement/Jury -------------------------------------------------------------
\begin{center}
\begin{tabular}{l l l l}
\multicolumn{2}{l}{\textbf{\textsc{Thèse} dirig\'ee par : }} & & \\[.5ex]
	& \textbf{Christophe \textsc{Guimbaud}} &  &  Professeur, Université d'Orléans\\
	& \textbf{Fatima \textsc{Laggoun}} & &  Directeur de Recherche, ISTO, Orléans\\[1ex]
\multicolumn{2}{l}{\textsc{\textbf{Rapporteurs :}}} & & \\[.5ex]
%\begin{tabular}{l p{1cm} p{10cm}}
	& \textbf{Denis \textsc{Loustau}} & &  Professeur, Université de Franche-Comté\\
	& \textbf{François \textsc{Gillet}} & &  Directeur de Recherche, INRA\\[1ex]
\hline\\ [-1ex]
\multicolumn{2}{l}{\textsc{\textbf{Jury : }}} & &\\[.5ex]
	&\textbf{Pr\'enom \textsc{Nom}} &  & Titre, \'etablissement, Pr\'esident du jury\\
	&\textbf{Pr\'enom \textsc{Nom}} &  &  Titre, \'etablissement\\
	&\textbf{Pr\'enom \textsc{Nom}} &  &  Titre, \'etablissement\\
	&\textbf{Pr\'enom \textsc{Nom}} &  &  Titre, \'etablissement\\
	&\textbf{Pr\'enom \textsc{Nom}} &  &  Titre, \'etablissement\\
	&\textbf{Pr\'enom \textsc{Nom}} &  &  Titre, \'etablissement\\
	&\textbf{Pr\'enom \textsc{Nom}} &  &  Titre, \'etablissement\\
\end{tabular}
\end{center}

\vfill

\restoregeometry % restore base margins
\usefont{T1}{cmr}{m}{n} % restore base font

\pagenumbering{roman} % On active la numérotation des pages

\pagestyle{frontmatter}% frontmatter page style

\tableofcontents
\addcontentsline{toc}{chapter}{Table des matières}
\listoffigures
\addcontentsline{toc}{chapter}{Liste des figures}
\listoftables
\addcontentsline{toc}{chapter}{Liste des tableaux}
% AVANT PROPOS

\chapter{Avant-propos}
%\addcontentsline{toc}{chapter}{Avant-propos}

% REMERCIEMENTS

\chapter{Remerciements}
%\addcontentsline{toc}{chapter}{Remerciements}


%%%%%%%%%%%%%%%%%%%%%%%%%%%%%%%%%%%%%%%%%%%%%%%%%%%%%%%%%%%%%%%%%%%%%%%%%%%%%%%%
% Mainmatter
%%%%%%%%%%%%%%%%%%%%%%%%%%%%%%%%%%%%%%%%%%%%%%%%%%%%%%%%%%%%%%%%%%%%%%%%%%%%%%%%
\mainmatter
\doublespacing
\pagestyle{mainmatter}% mainmatter page style

% INTRODUCTION
% OBJECTIFS : Introduire le contexte générale, l'échelle globale des problématiques
% Cette introduction doit être parfaitement compréhensible par un béotien !

\chapter*{Introduction}
\markboth{Introduction}{}
\addcontentsline{toc}{chapter}{Introduction}
\newpage
\begin{linenumbers}


\section*{Contexte général}

%PLAN : 
%- Le changement global, une thématique actuelle
%	* Historique du changement climatique ?
%	* Consensus scientifique ?
%	* Qu'est-il attendu  en terme de changement (réchauffement ?)
%- Le cycle du carbone : des flux et des stocks
%	* importance relative des stocks et flux
%	* principaux facteur forcant à l'échelle globale ?
%- Et les zones humides et les tourbières dans tout ça ?
%	* Qu'est ce qu'une tourbière ?
% 	* Fonctionnement vis à vis du cycle de carbone 
% 	* Une sensibilité particulière (localisation, anthropisation)
% 	* Écosystème non pris en compte dans les modèles globaux

En 1957, Charles David Keeling, scientifique américain, met au point et utilise pour la première fois, un analyseur de gaz infra-rouge pour mesurer la concentration de \COO de l'atmosphère sur l'île d'Hawaii, à Mauna Loa.
La précision et la fréquence importante de ses mesures lui permirent de voir pour la première fois les variations journalière et saisonnière des concentrations en \COO atmosphérique, mais également à plus long terme leur tendance haussière \cite{harris2010}.
Le \COO est un gaz à effet de serre (GES) et son accumulation dans l'atmosphère... \todo[inline]{force ? comparaison ? explication effet de serre ?}
Ce constat a probablement joué un rôle considérable dans la prise de conscience, par la communauté scientifique, de l'importance et de l'intérêt de l'étude du changement climatique et plus largement des changements globaux.
Car si à l'époque les concentration en \COO était inférieure à 320~ppm (partie par millions) elles ont dépassées, au printemps 2014, la barre symbolique des 400~ppm selon un communiqué de l'Organisation Météorologique Mondiale. Les concentrations pré-industrielles (avant 1800) sont quand à elles généralement estimée à 280 ppm \cite{Siegenthaler1987}.

Aujourd'hui, que ce soit pour le comprendre, le caractériser ou bien le prédire, de nombreux \todo{Combien ? cf fact sheet IPCC} scientifiques dans un grand nombre de disciplines, travaillent directement ou indirectement sur les changements globaux.
Ils sont nombreux également à collaborer au sein du  Groupe d'experts Intergouvernemental sur l'Évolution du Climat (GIEC), qui rassemble, évalue et synthétise les connaissances internationales liée au sujet.
%Des expérimentateurs, des modélisateurs, des climatologues, des atmosphéristes \todo{atmosphéristes, vérifier que ce mot existe...}, des écologues, qu'ils travaillent sur des archives ou des données actuelles. \todo{à reformuler}

De manière générale, parmi les flux de C mesurés entre la biosphère et l'atmosphère, la respiration et la photosynthèse sont les plus  important, 98 et 123 PgC/yr pour le flux de respiration globale (+ les feux) et la photosynthèse respectivement \cite{Bond-Lamberty2010,Beer2010}. Pour comparaison les flux liés à la production de ciment et aux ressources fossiles (charbon, pétrole et gaz) représentent 7.8 PgC/yr \cite{Ciais2014}.

Étroitement lié aux changements globaux, le cycle du carbone est particulièrement étudié, quels sont les réservoirs, quels sont les flux et comment vont-ils évoluer ? 
\todo[inline]{schéma ?}


Zones humides tourbières

historique des tourbières, généralités sur l'histoire des tourbières vis à vis des hommes
Sujets principaux qui ont menés à l’étude des tourbières jusqu'à nos jours (Exploitation, effet de serre)

Pourquoi étudier les tourbières aujourd’hui ? 





L'étude des tourbières se poursuit car, en plus de rendre de nombreux services écologiques \index{services ecologiques@services écologiques} (épuration du sol, régulation des flux hydriques, biodiversité), elles constituent un stock de carbone relativement important au regard de la surface qu'elle occupent. Ainsi il est généralement admis que les tourbières contiennent un quart à un tiers du carbone présent \todo{Chiffres (surfaces...)} dans l'ensemble des terres émergées tandis qu'elle ne constituent que 3 \% des surfaces continentales \plop. Ce ratio relativement important, correspond à un stock d'environ 455 Gt \cite{Gorham1991,Turunen2002} est à mettre perspective avec les autres stock du cycle du carbone. On observe que ce stock est du même ordre de grandeur que celui de la végétation 

%Les tourbières sont des écosystèmes qui, s'il ne s'étendent pas sur une surface très importante (leur surface est estimée à 3\% des surfaces émergées) contiennent, relativement à leur extension, une quantité de carbone importante (455 Gt selon \cite{Gorham1991}).
En conséquence dans un contexte **d'augmentation des GES dans l'atm et de réchauffement**, l'évolution de ce stock, sa pérennité ou sa remobilisation est un sujet d'étude important. De plus cette importance n'est à ce jour pas prise en compte de façon spécifique dans les modèles climatiques globaux.

En France les tourbières s'étendent sur environ 60 000 Ha (\plop).

\begin{table}[!h]
\centering
\begin{tabular}{lll}\toprule
Compartiment & Stock (en Gt de C) & référence \\ \midrule
Tourbières & 455 & \cite{Gorham1991} \\ 
Végétation & 450 -- 650 & \cite{Robert2003}\\ 
Sols & 1500 -- 2000 & \cite{Robert2003,Post1982,Eswaran1993}\\ 
\COO atmosphérique & 750 -- 800 & \cite{Robert2003}\\ 
Permafrost & 1700 & \\ \bottomrule
\hline 
\end{tabular}
\caption{Estimations des stocks de C pour différents environnements}
\end{table}

\todo{Pas d'entrée "journal" pour Post1982}


Transition modèles

En octobre 2013 le Groupe d'experts Intergouvernemental sur l'Évolution	du Climat (GIEC) a publié le rapport du groupe de travail I qui travaille sur les aspects scientifique physique du système et du changement climatique.
S'il note que les connaissance ont avancées, il note également que de nombreux processus ayant trait à la décomposition du carbone sont toujours absent des modèles notamment en ce qui concerne le carbone des zones humides boréales et tropicales et des tourbières. \plop



\section*{Objectif de la thèse et approche mise en oeuvre}

%Dans ce contexte l'objectif de ces travaux est de comprendre la dynamique des flux de carbone.
%Principalement le \COO qu'il soit gazeux, dissout ou particulaire et le \CHH.
%Notamment caractériser la variabilité spatiale et temporelle et déterminer quels sont les facteurs de contrôles dominants.

L'objectif de ces travaux est donc de mieux comprendre la dynamique du carbone au sein des tourbières.
Tout d'abord en caractérisant la variabilité spatiale et temporelle des flux de carbone à travers l'établissement de bilan de carbone.
De déterminer quels facteurs environnementaux contrôlent le fonctionnement comme puits ou source de carbone de ces écosystèmes.
Enfin construire, dans un esprit de synthèse et d'ouverture et à l'aide des connaissances acquise, un modèle intégrateur permettant un lien avec les modèles globaux et notamment ORCHIDE, afin que ces écosystèmes puissent être pris en compte à cette échelle.

Pour atteindre ces objectifs, nos travaux ont été articulés autour de trois volets \todo{volet... t'as pas mieux ? Branche ? -\_-"} principaux :
Dans un premier temps, l'\textbf{observation} régulière des flux de gaz (\COO et \CHH) ainsi que d'un certain nombre de paramètres environnementaux servant à la caractérisation des variabilités spatiales et temporelles, ainsi qu'à l'étude des facteurs contrôlant.
Certains facteurs contrôlant qui sont, dans un second temps, étudiés plus spécifiquement à travers un volet \textbf{expérimentation}.
Ce dernier doit permettre une meilleure compréhension de processus clé avec notamment l'impact de l'hydrologie.
Enfin un troisième volet axé sur la \textbf{modélisation}, avec le développement d'un modèle le plus mécaniste possible.


% Contenu des chapitres de la thèse .
Cette thèse est structurée de la façon suivante :
Le chapitre 1 est une synthèse bibliographique, un état de l'art des connaissances liées au sujet.
Les chapitres 2 et 3 rassemblent les travaux du volet observation, ils concernent respectivement le suivi XX et le suivi YY 
Les chapitres 4 et 5 développent la partie expérimentale à travers l'impact d'un assèchement et celui d'un rehaussement du niveau de l'eau.
Le chapitre 6 concerne plus spécifiquement la modélisation, même si ce volet interviendra par ailleurs de façon transverse dans les autres chapitres.
Enfin une conclusions et des perspectives seront exposées.

\end{linenumbers}

% % % % TRASH ? REUTILISATION APRES.
%
%Afin de répondre à notre problématique (La variation spatiale et temporelle des flux de \COO $\leftarrow$ Ceci est la thématique pas un problématique...), Trois approches complémentaires ont été mises en œuvres. 
%Tout d'abord la mise en place de suivis terrain, qui nous ont permit de mieux comprendre certains processus. 
%L'expérimentation ensuite, sur le terrain ou en laboratoire, afin d'étudier plus en détail l'impact d'un facteur forçant majeur qu'est l'hydrologie. 
%Et enfin la modélisation afin de synthétiser ces connaissances dans un cadre dont on espère pouvoir se resservir dans différents sites.

% % CHAPITRE 1 : SYNTHESE BIBLIO
%
%Réflexions
%
%NE PAS OUBLIER DE CALER UN MAXIMUM D'ORDRE DE GRANDEUR ! 
%	(émissions CO2, CH4, stocks flux, surface des tourbières, végétation ...)
%
%Bien différencier l'intro générale qui doit être lisible par un béotien, de l'intro au travail de la thèse qui doit être un état de l'art précis et documenté sur les travaux antérieurs (synthèse biblio)
%
%Ou caler la partie de biblio sur l'expérimentation ... (peut être dans la synthèse biblio, paragraphe "approche mise en oeuvre"
%INTRODUCTION GENERALE (à mettre dans le chapitre Intro)
%
%- Qu'est ce qu'une tourbière ? (Éventuellememt comment se forme-t-elle ?)
%	*Définition
%	*Formation/Évolution (stockage du C, battement de la nappe ?)
%	*Classification
%- Les tourbières et les hommes 
%	*Usages d'hier et d'aujourd'hui (Combustible, horticulture, matériau de construction)
%	*Les thématiques scientifiques (pourquoi les avoir étudier et les étudier en gros)
%	*Le contexte dans lequel va s'inscrire le travail qui suit
%
%SYNTHESE BIBLIOGRAPHIQUE
%
%- Quelles sont les grandes thématiques de recherche liées aux tourbières ?
%	*Exploitation
%	*Archives
%	*Émissions de GES
%- Plus précisément quelles sont les grands axes de recherche sur ces écosystèmes et liés aux émissions de GES.
%	*Processus de création de GES (CO2 et CH4) (Facteurs contrôlant généralement invoqués)
%	*Processus de migration des GES dans le profil
%	*Processus de stockage/capture
%- Approches mise en oeuvre
%	*Modélisation (empirique et mécaniste)
%	*Expérimentation (différentes techniques pour mesurer les émissions de GES, différentes techniques de chambre...)
%	*Variabilité spatio-temporelle (notion d'échelle)
%	
%	DOIS-JE TRAITER
%	La classification des tourbières ?
%	Hummock and Hollow ? Dire qu'on n'est pas dans ce niveau de détail ?
%	
%
%HISTORIQUE (études concernant les tourbières)
%
%1968-1969  Boelter : propriétés physique des tourbes
%1977 Boelter : hydrologie, caractéristique des sols organiques, chimie des écoulements
%1978 Ingram : Classification
%1981 Parkinson : Déjà l'amélioration d'une méthode pour la mesure des émissions de la respiration du sol
%1984 Clymo : Les limites à la croissance des tourbières
%1986 Chason : Conductivités hydraulique et propriétés physiques
%1989 Moore : Influence du niveau de la nappe sur les émissions de CO2 et de CH4
%//1990 Raich : Comparaison de 2 méthodes de chambre statique pour mesurer les flux de CO2
%1991 Gorham : Rôle des tourbières dans le cycle du carbon et réponse au changement climatique
%//1992 Raich : Flux de CO2 dans la respiration du sol et relation vis à vis du climat et de la végétation
%1992 Roulet : Flux de méthane (fen) et changement climatique
%1993 Bubier : Émission de méthane dans les zones humides
%1993 Bubier : Microtopographie et flux de méthane dans tourbières boréales.
%1993 Abbès : Sorption de l'ammonium ? (ammonia) par la tourbe et fractionnement de l'azote
%1994 Lloyd : Dépendance de la respiration du sol à la température
%1994 Bubier : Perspective écologique sur les émissions de méthane dans les zones humide de l'hémisphère nord
%//1994 Nay : Biais des méthodes de chambre pour la mesure des flux de CO2
%1995 Kirschbaum : Dépendance à la température de la décomposition de la MO (effet sur stock de C et changement Clim)
%1995 Bubier : Prédiction des émissions de méthane à partir de la distribution des bryophytes (tourbière hémisphère nord)
%1995 Bubier : Contrôles "écologique" sur les émissions de méthane dans les tourbières de l'hémisphère nord
%1995 Bubier : Relation entre la végétation avec les émissions de méthane et les gradients hydrochimique
%//1995 Bekku : Mesure de la respiration du sol avec une méthode de chambre fermée (IRGA)

% PLAN (2015-03-03)
%I. Définitions
%1 Tourbières/Tourbe (surface, type, localisation, biodiversité, services écologiques...)
%2 Classification
%3 Historique
%	a Utilisation
%	b Études scientifiques
%	
%Transition : Réaction aux changements globaux (comment fonctionnent-elles ?)
%
%II. Fonctionnement
%1 Stock
%2 Flux
%	a Entrants (Photosynthèse)
%	b Sortants (Méthanogénèse, Respiration)
%3 Facteurs Contrôlant
%	a Hydrologie (WTL,HR)
%	b Propriétés physiques (T, densités, conductivités thermiques...)
%	c Végétation (Bryophytes/Vasculaires)
%	d Météo

\chapter{Synth\`{e}se Bibliographique}
\newpage

Dans ce chapitre, nous commenceront par donner une vue de ce que sont les tourbières : Que sont-elles ? Depuis quand sont-elles étudiées ? Pourquoi les a-t-on étudiés ?
Nous continuerons en entrant plus en détails sur leur fonctionnement vis à vis des flux de carbone.
Enfin nous verrons quels sont les facteurs contrôlant majeurs de ces flux.
\section{Qu'est ce qu'une tourbière ?}
Les tourbières font partie d'un ensemble d'écosystèmes que l'on appelle les zones humides.
Ces écosystèmes ont, comme leur nom peut le laisser penser, un lien étroit avec l'eau.
Ainsi outre les tourbières, les lacs, les deltas, les mangroves en font également partie.
Les tourbières représentent 50 à 70 \% des zones humides (francez2000, wise use of peatland), s'étendant sur plus de 4000000 de km$^{2}$ (Lappalainen).
Une définition régulièrement utilisée pour caractériser ce qu'est une tourbière est : "Tout écosystème possédant au moins 30 cm de tourbe".
Malgrès le fait que cette définition ne soit pas très éclairante, elle est souvent préférée à une définition plus stricte : "écosystème dans lequel un processus de tourbification est actif". qui n'est pas beaucoup plus lumineuse.
Qu'est ce que la tourbe ?
Et qu'est ce que la tourbification ?
La tourbe est le résultat de l'accumulation et de la, faible, dégradation de litières végétales.
C'est ce que l'on appelle la tourbification.
La structure de la tourbe est liée à une production primaire particulièrement importante, comparée à un sol forestier par exemple, mais plutôt à une dégradation plus faible des matières organiques.
\textbf{Attention tu parles ici des tourbières boréales spécifiquement sans les avoir décrites auparavant.}
Cette faible dégradation est due notamment au fait que ces écosystèmes, sont saturés en eau une grande partie de l'année (ce sont des zones humides). 
Par conséquent, l'accès à l'oxygène est plus difficile diminuant d'autant l'activité aérobie, dont la respiration des micro-organismes et des plantes.

%Que sont-elles ? 
%Les tourbières font partie d'un ensemble d'écosystèmes que l'on appelle les zones humides.
%Dans ces zones humides, en plus des tourbières sont classés des écosystèmes comme les lacs, les deltas, les mangroves.
%Elles sont définies comme les écosystèmes possédant au moins 30 cm de tourbe.
%Cette règle peut cependant changer selon les régions.
La classification des tourbières n'est pas simple, de nombreux critères existent selon leur mode de formation, leur source d'eau, leur physico-chimie.
Au niveau terminologie, de nombreux termes ont été utilisés parfois en contradiction les uns avec les autres.
Les tourbières sont donc des zones dans lequelles l'eau à une rôle majeur.
%Notamment le niveau de l'eau étant, normalement, élevé il empêche la décomposition des litières en diminuant fortement l'accès à l'oxygène et donc la respiration.
Les tourbières sont le siège d'une biodiversité spécifique relativement importante et rendent un certain nombre de services écologiques.
Parmi la végétation caractéristique de ces écosystèmes, les sphaignes, des bryophytes (des mousses) sont normalement présentes en abondance.Ces écosystèmes on été et sont encore perturbés par différentes activités humaines, notamment l'agriculture, l'utilisation de la tourbe comme combustible, et comme substrat horticole.
Ces perturbations peuvent induire des modifications de fonctionnement, notamment l'envahissement de ces écosystèmes par une végétation vasculaire.
Ces écosystèmes s'étendent sur 4 000 000 de km2.

Depuis quand sont-elles étudiées
D'abord étudiées pour leurs propriétés physiques afin de connaitre leur qualité en tant que combustible.
Elle sont maintenant majoritairement étudiée à travers le prisme des changements globaux.
Ainsi les études concernent les flux de GES, ...

\subsection{Le CO2 processus de productions/dégradations et de transport}
Le CO2 est un des principaux gaz à effet de serre si bien que les autres sont souvent classés en fonction de ce dernier.

** Historique précis des études concernant les GES (CH4)

\subsubsection{La production de CO2 : La respiration de l'écosystème (processus de production)}
Dans les tourbières le CO2 est produit par des sources multiples.
Ces sources sont la respiration des de la flore qu'elle soit aérienne ou souterraine et la respiration microbienne.
Une autre source de CO2 est l'oxydation du CH4 lors de sa migration des zones anoxiques aux zones oxiques de la colonne de tourbe.
Enfin dans les zones anaérobie, le CO2 peut être produit par fermentation (respiration anaérobie).
La production de CO2 est donc un signal intégré sur l'ensemble de la colonne de tourbe. 
C'est cette multitude de processus qui rend l'estimation de ce flux difficile, en effet chacune des respirations n'aura pas la même sensibilité vis à vis de facteurs contrôlant.
La respiration de l'écosystème (RE) est définie comme l'ensemble des respirations de la colonne de tourbe, en incluant à la fois sa partie aérienne et sa partie souterraine.
La respiration du sol (SR) est elle définie comme l'ensemble des respirations de la colonne de tourbe, en excluant la partie aérienne.
La respiration du sol comprend donc principalement les respirations issues de la rhizosphère et des communautés de micro-organisme.

Les tourbières sont des écosystèmes dont la production primaire est estimée à environ 500~g~C.m$^{-2}$ \cite{francez2000}. 
La strate muscinale pouvant jouer/participer/produire jusqu'à 80\% de la production primaire \cite{francez2000}.
Cette production primaire n'est pas particulière élevée \plop et c'est en fait la faible décomposition des matières organiques qui permet aux tourbières de stocker du carbone.
L'accumulation moyenne estimée dans les tourbières boréales est de 30~g~C.m$^{-2}$. Le taux d'accumulation varie en fonction des espèces végétales présentes (\plop), le niveau d'eau (\plop), ... (??)

\subsubsection{La capture du CO2 :(processus de stockage)}
C'est évidemment par photosynthèse que le CO2 est pompé de l'atmosphère pour être stocké dans tissus des végétaux avant d'être en partie non dégradé et donc stockés dans les litières puis dans la tourbe à proprement parler
La vitesse de stockage a pu varier au cours du temps mais elle est estimé à XXXX, ainsi la majorité des tourbières actuelles ont un stock qui remonte à quelques milliers d'années.
Les estimations précise du stock de C présent dans ces écosystèmes sont délicates, à la fois car la définition de ce qu'est une tourbière que varier selon les régions, mais également car leur étendue exacte n'est pas triviale à estimer, pas davantage que leur profondeur moyenne.
Cependant il est usuellement admis que le stock de carbone se situe entre 270 et 500 Gt de C

\subsection{Le CH4 processus de productions/dégradations et de transport}
Comparé au CO2, le CH4 est un GES qui est bien moins présent dans l'atmosphère (CHIFFRES!).
Cependant son "pouvoir de réchauffement" est bien plus important (effet radiatif CO2 x 100) (CHIFFRES !) (D'abord la vapeur d'eau, ensuite le CO2 et enfin le CH4)
Il est usuellement convenu (???? ref) que dans une tourbière le méthane représente environ 5\% du bilan de C.

\section{Les facteurs majeurs contrôlant les flux}
\subsection{La température et les flux}
La température est le premier facteur contrôlant les flux.
Comme pour toute (la majorité ? y a-t-il des réactions chimiques non influencées par la température ?) les réactions chimique la température influe sur les vitesses de réactions. 
Plus la température augmente plus la vitesse de réaction augmente.
La température à donc un rôle important à jouer au niveau des flux.

température et ER
Concernant la, ou plutôt les respirations de l'écosystème, l'influence de la température sera différente selon les sources considérés.
Ainsi à la fois les plantes et les communauté de micro-organisme ne réagiront probablement pas de la même façon, au mêmes moments et avec les même intensités.

température et NEE
Pour la NEE même s'il semble y avoir moins de sources possible puisque seule la végétation photosynthétique est concernée, l'influence de la température est également fonction de la végétation présente.
De plus ce signal est plus ou moins covariant avec la luminosité ce qui ne facilite pas son interprétation.

Synchronisation


\subsection{L'hydrologie dans les tourbières et l'effet sur les flux}
L'hydrologie est comme nous l'avons précisé un peu plus haut, un facteur d'une grande importance dans les tourbières.
Nous distinguerons ici le niveau de la nappe qui est la hauteur sous la surface du sol permettant d'accéder à la zone saturée ? à l'eau "lirbe" ?
Et l'humidité du sol qui est une estimation de la quantité d'eau présente dans la zone non-saturée.

\subsubsection{L'effet du niveau de la nappe}
Le niveau de la nappe est important car il sépare la colonne de tourbe en une zone oxique, ou il y a présence d'oxygène, et une zone anoxique dans laquelle l'oxygène est absent.
Ces deux zones vont avoir des comportements différents.
La zone anoxique, sous le niveau de la nappe, est une zone dans laquelle la production de CO2 est très faible car sans oxygène seule les processus de respiration anaérobie peuvent avoir lieu.
Par contre dans c'est dans cette zone que sera produit le méthane.
La zone oxique, proche de la surface, va permettre à la fois aux racines et aux micro-organismes de respirer.
Cette zone est donc l'endroit ou est produit la majorité du CO2, l'endroit ou la matière organique est le plus dégradée.
Lors de la migration du méthane dans la colonne de tourbe ce dernier aura tendance à être oxydé en CO2 lors de son passage dans cette zone oxique.
Certaines plantes permettent cependant au méthane de passer à travers l'aerenchyme et d'éviter ainsi d'être oxydé.

\subsubsection{L'effet de l'humidité relative}

\subsubsection{Résilience de la tourbe}
Les propriétés physique de la tourbe jouent bien évidemment un rôle important sur cette capacité de rétention d'eau.
Cependant dans le cas d'épisode de sécheresse important, il a été constaté que ces capacités n'était pas immédiatement recouverte en totalité.


\subsection{La végétation dans les tourbières et l'effet sur les flux}
Les communautés végétales évoluent en parallèle de l'évolution de la tourbière (succession végétale).
Les tourbières sont le siège d'une végétation caractéristique : Les sphaignes.
Ces bryophytes sont la clef de voûte de ces écosystèmes d'abord parce que leur litière sont moins facilement dégradable que celle des espèces vasculaires.
Ensuite parce qu'elle favorisent dans leur environnement local, les conditions favorable à leur développement. 
On les appelle d'ailleurs des espèces ingénieures.
Ces végétaux sans racines ont également une grande capacité à retenir l'eau (ce sont de véritables éponges) retenant également les nutriments. 
Ceci favorisant un milieu pauvre en nutriment et donc défavorable aux autres espèces (vasculaires?).
Il existe un grand nombre d'espèce de sphaignes (CHIFFRES+REF).
Par la suite il ne sera pas fait de distinction entre les différentes espèces présentes sur les différents sites étudiés.
Cependant dans de nombreuses tourbières on constate un envahissement par des végétaux vasculaires.
Ces plantes, sont souvent des pins, des bouleaux et des molinie ?
Elles ont un effet sur la production de CO2 principalement en aérant le sol, permettant à l'oxygène de migrer plus loin dans le profil, permettant à l'activité aérobie (plus efficace) d'agir sur une plus grande profondeur.
Ces végétaux peuvent également pomper de l'eau en quantité (arbre?) ?


% % % % % % % % % % % % % % %

\todo[inline]{Descriptif et comparaison des méthodes permettant de mesurer les flux de gaz}



%\part{Observation}
% CHAPITRE 2
% SUIVI VARIABILITE SPATIALE

\chapter{Sites d'études et méthodologies employées}
\newpage

\section{Présentation des sites d'études}

L'ensemble des sites d'études sont regroupés au sein d'un service d'observation

\subsection{La Guette}

La tourbière de La Guette est situé à Neuvy-sur-Barangeon, en Sologne, dans le département du Cher.
Le site s'étend sur une surface d'une vingtaine d'hectare avec une géométrie relativement allongée.
Avec une conductivité généralement inférieur à 80 uS/m2 et un pH compris entre 4 et 5 elle se classe parmis les "transitionnal poor fen"
Les datations effectuées sur le site permettent de dire que la tourbière est agée de 5 à 6000 ans.
Dans les années 19XX la construction d'une route coupe la tourbière dans sa partie sud.
En 2008 le récurage du fossé de drainage bordant la route semble entrainer une augmentation significative des pertes d'eau du système.

Des travaux (SOURCE, Émelie) d'analyse de photos aériennes ont ainsi montré une progression importante du boisement (principalement des pins (pinus Sylvestris) et des bouleau (Betula sp.). Des herbacées envahissent également le site avec une forte présence de la molinie (Molinia caerulea)

Sont présente sur le site un certain nombre d'espèces caractéristiques des tourbières comme les sphaignes (principalement Sphagnum cuspidatum et Sphagnum rubellum) et Eriophorum augustifolium.


\subsection{Frasne}

\subsection{Landemarais}

\subsection{Bernadouze}


Au sein de ses sites de nombreuses mesures ont été effectuée et notamment des mesures de flux de GES à la fois concernant le CO2 et le CH4. La méthodologie étant transverse à de nombreuses expérimentations il convient de l'expliquer au préalable.

\section{Mesures de flux}

Il existe de nombreuses façon de mesurer des flux de gaz. 
Des méthodes globale comme les tour à flux utilisant des méthodes d'Eddy Covariance. 
Des méthodes plus locale, les chambres d'accumulation de gaz qui peuvent être statique ou dynamique, selon que la sonde mesurant le gaz soit directement dans la chambre ou que le gaz soit apportée à cette dernière via un système de pompe. Elles peuvent être ouverte ou fermée.

La méthode de mesure retenue pour ces travaux est l'utilisation de chambre statique fermée, permettant une mesure locale et directe des flux.
Pour cela des embases sont placées sur le terrain. Il s'agit de cylindre de PVC d'une dizaine de cm, percés dans leur partie basse afin de minimiser les impacts sur les flux d'eau et sur de développement racinaire et enfoncé dans le sol.
Les embases sont généralement posée 12h avant toute mesure afin de ne pas mesurer de dégagement gazeux liés à l'installation.

Que mesure-t-on ?
Le plus souvent 2 mesures consécutives sont effectuées la première avec une chambre transparente permettant d'accéder à la NEE et l'autre avec une chambre recouverte d'un isolant permettant de bloquer la lumière et permettant de mesurer les respirations. (pourquoi les respirations?)

De nombreux écueils peuvent rendre une mesure inexploitable. D'abord le placement de la chambre, cela peut sembler trivial mais positionner la chambre au milieu d'herbacées et de bruyère n'est pas tourjours évident. Plus anectdotiquement des sphaignes gelées, recouvrant les bords de l'embase rendent la pose de la chambre difficile voire impossible. Selon l'heure de la journée des gradients de concentrations peuvent être présent et augmenter localement les concentrations de CO2 de façon importante allant jusqu'à saturer la sonde.

QUESTIONS :

*Taille des embases ? Effets de bord ?
*Perturbation du milieu ? (Mesure de végétation, pose de la chambre, mesure pièzo...)
*Impact de la strate arborée ?
*Validité des profils de température ?
Méthode de Chambre fermée (Biais ?)

Améliorations ? \todo{Lister les amélioration à faire ou non}


\section{Facteurs contrôlants et suivi des flux}
Afin de déterminer l'impact de facteurs contrôlants sur ces flux, mesurer les flux ne suffit pas il faut également mesurer les variables environnementales dont on pense qu'elles seront des facteurs contrôlants important. Le nombre de ces variables et la méthodologie employée pour les mesurer étant différente selon l'expérimentation ils seront détaillés au fil des besoins par la suite.

% CHAPITRE 3
\chapter{Effets de la températures sur les variations journalière des flux de CO2}
\newpage

\section{Introduction}
Les flux de gaz et notamment les flux de CO2 sont fonctions de la température.
La température dépend quand à elle de l'énergie reçue par le soleil et donc varie de façon journalière, saisonnière et au delà !

Afin de palier à ces deux aspects un autre suivi a été mis en place : l'étude des flux de \COO à relativement haute fréquence \todo[inline]{combien ? qu'est ce qu'une haute fréquence ?} pendant 3 jours et sur 4 sites différents \todo[inline]{liste des sites ?}.

Nous avons donc avec ces deux suivis, une vision à la fois sur la variabilité spatiale, au sein d'un site ou inter-site, et une vision sur la variabilité temporelle quelle soit saisonnière, annuelle ou journalière.

Ce schéma n'est bien sur pas parfait, ainsi les sites étudiés restent des sites situés en France alors que la majorité des tourbières se situent à des latitudes plus élevées, dans les zones boréales et sub-boréale.

\todo[inline]{Proportion des tourbières qui ont été exploités ? qui sont encore à l'état naturel ? à mettre en regard avec la représentativité d'une tourbière comme La Guette. Est-elle représentative ? La majorité des tourbières sont perturbées... Sont-elles envahies par des végétaux vasculaires ?}

L'étude d'un système complexe de façon globale permet d'avoir une vision globale, cependant il est difficile de comprendre certains processus quand s'ils sont noyés dans un tel système. 
L'expérimentation, qu'elle soit sur le terrain ou en laboratoire permet de simplifier notre système afin de pouvoir déterminer l'impact de tel ou tel facteur plus particulièrement, afin de mieux comprendre tel ou tel processus.
Ainsi ont été mis en place différentes expérimentation bla bla bla.

\section{Présentation de l'expérimentation}

La respiration de l'écosytème (Re) est mesurée tous les quarts d'heure avec une méthode de chambre fermée.
La chambre, en plexiglas, est recouverte d'un isolant, un ventilateur placé à l'intérieur de la chambre permet d'homogénéiser l'air.
Ce dernier permet d'oculter la lumière du jour, et de conserver une température à l'intérieur de la chambre proche de la température extérieure.
Le \COO est mesuré à l'aide d'une sonde Vaisala (\plop précise).
Chaque mesure dure au maximum 5 minutes, délai permettant d'avoir une stabilisation du flux après la pose de la chambre et suffisant de points pour avoir une pente claire.

Les mesures sont faites en continu pendant 72h sur 4 embases. 
Chaque embase est donc mesuré une fois par heure et l'ordre des mesures a été déterminé de façon aléatoire.

En plus des mesures de \COO un piézomètre et une station météo a été installé à proximité des embases.
La station météo nous permet d'aquérir des données à haute fréquence (1 Hz, une mesure par seconde).
Les paramètres suivis sont, la radiation solaire, la température de l'air à 5 cm, la température du sol à différentes profondeurs (5, 10, 20, 30 cm) et l'humidité.

Des profils de températures réalisés (avec quelle sonde ?) ponctuellement dans les embases permettent de recaler chaque embase par rapport aux profils de la station.

Des mesures de NEE ont été testée, la première série sur la tourbière de LaGuette en utilisant le protocole de la variabilité spatiale (à préciser)
LE problème de ce protocole est l'augmentation de la température à l'intérieur de la chambre.
Cette augmentation peut engendrer dans les cas extrêmes une différence de température de plus de 10\degres C et entrainer l'arrêt de la photosynthèse dans la chambre.
(Probablement par fermeture des stomates des végétaux.)
Pour pallier à ce problème des "bloc de froid" ont été utilisé afin de minimiser la différence de température entre l'air à l'intérieur et à l'extérieur de la chambre.
Cette solution permet de diminuer la différence de température, mais il est difficile de contrôler précisément la température...
Un autre souci lors de l'expérimentation a été la perturbation de la végétation.
Répéter aussi régulièrement les mesures pertube la végétation sur 4 à 5 cm de part et d'autre de l'embase.

\section{synchronisation et profiles (article)}

%\part{Expérimentation}
% CHAPITRE 4
\chapter{Effets de l'hydrologie sur les flux de CO2 et CH4}
\newpage

\section{Introduction}
\section{Manipulation du niveau de l'eau en mésocosmes}
\subsection{Présentation de l'expérimentation}
\subsection{Résultats}

\section{Manipulation du niveau de l'eau (teneur en eau) in-situ}
L'étude des effets de l'hydrologie sur les émissions de flux de GES a également pu être menée directement in-situ au sein du projet CARBIODIV (Restauration hydrologique de la tourbière de La Guette : effets sur l'évolution de la biodiversité et le stockage du carbone.) dont l'objectif est de restaurer le fonctionnement hydrologique de la tourbière de La Guette.
\subsection{Présentation de l'expérimentation}
\subsubsection{Les travaux}
\subsubsection{Les stations scientifiques}
Deux stations ont été installées sur le site, dans deux sous-hydrosystèmes différents. Le premier en amont n'étant pas impacté par les travaux permet de contrôler les effets de site, et le second, en aval, enregistrera les effets de la restauration hydrologique.
\subsection{Résultats}

% CHAPITRE 5
\chapter{Effets de la végétation sur les flux}
\newpage

\section{Introduction}
\section{Mise en place d'un protocole}
\section{Impact des mesures de CO2 sur la végétation}
destruction...
% CHAPITRE 6
\chapter{Caractérisation de la variabilité spatio-temporelle des flux sur la tourbière de La Guette (Bilan de C)}
\newpage

\section{Introduction}
\section{Présentation du suivi}
Étudier la variabilité spatiale et temporelle des flux de carbone nécessite de mettre en place une observation régulière, un suivi adapté.
Qu'est ce qu'un suivi adapté ? Dans le cas des flux des gaz, on ne peut heureusement ou malheureusement pas suivre en permanence un site avec une fréquence importante (Qu'est ce qu'une fréquence importante).
Il faut donc choisir qu'elle(s) échelle(s) l'on souhaite étudier.
Lors de cette thèse nous avons choisi de mettre en place deux suivi différents et complémentaire.

Le premier est un suivi mensuel à l'échelle d'une tourbière. 
Plus précisément, les flux de \COO ont été mesurés, à minima, une fois par mois sur 20 points distribués sur l'étendue du site (La Tourbière de La Guette, Neuvy sur Barangeon, Cher, France).
À cause de contrainte technique et de leur importance relative \todo[inline]{Expliquer ici ou ailleurs que les flux de CH4 ne représente a priori que 5 \% du bilan de C sur une tourbière} les flux de \CHH ont été mesuré plus ponctuellement.
Ce suivi permet d'étudier la variabilité spatiale des flux au sein du site et cela sur une échelle temporelle de l'ordre de l'année qui permet d'avoir une vision sur les variations saisonnières.

Ce protocole d'observation seul ne nous permet cependant pas de rendre compte d'autres échelles de variabilité qui semble d'importance non négligeable. 
Tout d'abord la variabilité des flux journalière : 
Le système étant très lié à la photosynthèse les flux sont très dépendant de l’ensoleillement et donc des alternances jour/nuit
Ensuite qu'elle représentativité un site en particulier peut-il avoir par rapport à l'ensemble des sites existants ?

\subsection{Suivi des GES}
\subsection{Suivi des facteurs contrôlants}
Par ailleurs pour chaque embase le PAR a été mesuré, ainsi que la pression atmosphérique, la température de l'air et de la tourbe à différente profondeur (profils). Des mesures d'humidité ont également été effectué, 5 par embase afin de prendre en compte l'hétérogénéité présente.
\subsection{Suivi des flux liquides (DOC, POC)}
Mesures de DOC
Principe
Non Purgeable Organic Carbon
Pyrolyse oxydative
Ne mesure pas le CO2 dissous
L'acidification et le bullage permet d'expulser le CO2 dissous
Tout le carbone est ensuite pyrolysé le carbone est transformé en CO2 et mesuré
%\part{Modélisation}
% CHAPITRE 6
\chapter{Apport à la modélisation globale}
\newpage


\section{Introduction}
Enfin pour synthétiser l'ensemble de ces résultats, essayer d'en faire un tout cohérent et donc de comprendre notre système, nous l'avons modélisé en ayant pour objectif cet esprit de synthèse et de compréhension.
\section{Le modèle de Walter}

%%%%%%%%%%%%%%%%%%%%%%%%%%%%%%%%%%%%%%%%%%%%%%%%%%%%%%%%%%%%%%%%%%%%%%%%%%%%%%%%
% Backmatter
%%%%%%%%%%%%%%%%%%%%%%%%%%%%%%%%%%%%%%%%%%%%%%%%%%%%%%%%%%%%%%%%%%%%%%%%%%%%%%%%
\backmatter
\pagestyle{backmatter}% backmatter page style

% Conclusions ------------------------------------------------------------------
% CONCLUSIONS
\chapter{Conclusions et perspectives}
\newpage

Synthèse générale et discussion


Variabilité temporelle

Variabilité spatiale



% Bibliographie ----------------------------------------------------------------
\bibliography{manuscript}{}
\bibliographystyle{apalike}
\addcontentsline{toc}{chapter}{Références bibliographiques}
\clearpage

% Index ------------------------------------------------------------------------
\phantomsection % Permet "d'accrocher" hyperref
\addcontentsline{toc}{chapter}{Index}
\printindex

% 4e de couverture -------------------------------------------------------------
\thispagestyle{plain}
\singlespacing
%%%%%%%%%%%%%%%%%%%%%%%%%%%%%%%%%%%%%%%%%%%%%%%%%%%%%%%%%%%%%%%%%%%%%%%%%%%%%%%%
% 4e de couverture
%%%%%%%%%%%%%%%%%%%%%%%%%%%%%%%%%%%%%%%%%%%%%%%%%%%%%%%%%%%%%%%%%%%%%%%%%%%%%%%%
% set-up -----------------------------------------------------------------------
\pagenumbering{gobble} % stop page numbering
\newgeometry{left=2cm,top=1.5cm,right=1.5cm,bottom=1.5cm} % set up margin
\usefont{T1}{phv}{m}{n} % Set up font style
{\parindent0pt % disables indentation for all the text between { and }

% Résumé ---------------------------------------------------------------------
\begin{center}
		\large{Benoît D'ANGELO}\\ 
		\textbf{Variabilité spatio-temporelle des émissions de GES dans une tourbière à Sphaignes : effets sur le bilan de carbone}
\end{center}

\begin{framed}
	\begin{minipage}{\dimexpr\textwidth-2\fboxrule-2\fboxsep}
	Les tourbières ne représentent que 2 à 3\% des terres émergées mais stockent entre 10 et 25\% du carbone accumulé dans les sols.
	Les conditions de saturation en eau importante de ces écosystèmes diminuent la décomposition des matières organiques et favorise leur préservation.
	Ces écosystèmes sont cependant soumis à des contraintes anthropiques et climatiques importantes qui posent la question de leur devenir ainsi que celui du stock de carbone qu'ils hébergent.
	Une meilleure compréhension de ces écosystèmes est nécessaire afin de déterminer quels sont les facteurs qui contrôlent leur émissions de gaz à effets de serre (GES) et surtout comment les contrôlent-ils.\par
	Ce travail a donc consisté à suivre les émissions de GES et les facteurs contrôlant d'une tourbière de Sologne (La tourbière de La Guette) afin d'établir son bilan de carbone.
	En parallèle des expérimentations sur l'hydrologie ont été menées afin d'en préciser les effets sur les flux, et un suivi ponctuel sur différents sites a été réalisé dans le but d'étudier la variabilité à l'échelle journalière.\par
	Les résultats de ces travaux montrent que la tourbière de La Guette fonctionne comme un puits de carbone et ce malgré un niveau de nappe élevé ce qui suggèrent un effet de l'histoire antérieure du site.
	Il montrent également l'importance de la variabilité spatiale des flux que l'on peut estimer à l'échelle d'un site.
	Les expérimentations confirment l'importance de l'hydrologie et mettent en avant l'importance à haut niveau de nappe d'eau de phénomènes liés au transport des gaz entre leurs zones de production et l'atmosphère.
	Enfin l'étude de la variabilité journalière montre que la sensibilité à la température de la respiration peut être différente le jour et la nuit et que la synchronisation entre les températures du sol et la respiration peuvent améliorer la représentation de cette dernière.
Mots cl\'es : Tourbière, Gaz à effet de serre, \coo, \chh, bilan de carbone
	\end{minipage}
\end{framed}

\vfill

% Abstract ---------------------------------------------------------------------
\begin{center}
	\large \textbf{Spatio-temporal variability of Greenhouse gases emissions in a Sphagnum peatland: effects on carbon balance}
\end{center}

\begin{framed}
	\begin{minipage}{\dimexpr\textwidth-2\fboxrule-2\fboxsep}
	R\'esum\'e : (1700 caract\`eres max.)\par
	Lorem ipsum dolor sit amet, consectetur adipiscing elit. Proin volutpat ipsum id purus ultrices lobortis. Maecenas ornare enim quis eros. Nunc eget mauris ut quam malesuada mattis. Vestibulum ante ipsum primis in faucibus orci luctus et ultrices posuere cubilia Curae; Integer vel tellus. Nam rutrum, purus non sodales rhoncus, quam magna imperdiet eros, sit amet euismod justo metus at orci. Suspendisse neque turpis, feugiat interdum, faucibus vel, aliquet quis, risus. Etiam est elit, eleifend a, consequat sit amet, scelerisque nec, odio. Quisque id odio quis libero iaculis tincidunt. Sed non mi. Morbi aliquam commodo nibh. Integer justo purus, pulvinar a, suscipit vel, iaculis a, justo. Morbi ut orci. Maecenas fringilla orci. Phasellus auctor, enim vitae tempus egestas, justo mi cursus sem, vel blandit leo turpis vitae quam. Etiam sit amet felis vitae eros ornare porttitor.\par
	Curabitur felis velit, aliquam at, aliquet in, iaculis vitae, velit. Nunc lobortis magna id ligula. Vestibulum ante ipsum primis in faucibus orci luctus et ultrices posuere cubilia Curae; Integer congue ultrices mi.
	Isdem diebus Apollinaris Domitiani gener, paulo ante agens palatii Caesaris curam, ad Mesopotamiam missus a socero per militares numeros immodice scrutabatur, an quaedam altiora meditantis iam Galli secreta susceperint scripta, qui conpertis Antiochiae gestis per minorem Armeniam lapsus Constantinopolim petit.\par
Mots cl\'es : mot 1, mot 2, ...
	\end{minipage}
\end{framed}

% logos and university name ----------------------------------------------------
\includegraphics[width=0.2\textwidth, valign=c]{./images/logos/pucvl}
\hfill
%{\LARGE\textbf{UNIVERSITÉ D'ORLÉANS}}
\begin{minipage}{.5\textwidth}
\begin{center}
LPC2E/CNRS\\
3A, Avenue de la Recherche Scientifique\\
45071 Orléans cedex 2\\
France \\
\end{center}
\end{minipage}
\hfill
\includegraphics[width=0.2\textwidth, valign=c]{./images/logos/univ}

}

\listoftodos
\end{document}
